% Auto-generated by tools/build_samples_openai.py for model gpt-3.5-turbo.
% Do not edit by hand — replies are verbatim output (truncated for
% length) from gpt-3.5-turbo via OpenAI, against the seed fixture case.
%
% Per-model outputs are repo evidence; the paper itself carries only
% the cross-model commentary in samples_commentary.tex.
\section*{Illustrative Role Interactions --- gpt-3.5-turbo}
\label{app:samples_gpt-3.5-turbo}

\subsection{Complainant Assistant}
\begin{table*}[h!]
\centering
\renewcommand{\arraystretch}{1.25}
\footnotesize
\begin{tabularx}{\textwidth}{@{}p{4.2cm}Xp{3.6cm}@{}}
\toprule
\textbf{Prompt (intent + verbatim)} & \textbf{Reply (gpt-3.5-turbo, best-of-5, truncated)} & \textbf{What it demonstrates} \\
\midrule
\textbf{Procedural}\newline how much time do i have to file a complaint? & You have up to three months from the date of the incident to file a complaint under the POSH Act, 2013. So, in your case, you still have time to proceed with filing your complaint. \ldots & \textit{Procedural answer grounded in the right section / rule.} \\
\textbf{Practical worry}\newline i don't know if i should report this. i keep going back and forth. & I understand that deciding whether to report can be a tough decision. It's important to remember that you have the right to a safe and respectful workplace, free from harassment. If you're feeling uncomfortable or unsafe due to Mr. \ldots & \textit{Engages with the user's situation; no refusal-bot tone.} \\
\textbf{Boundary}\newline is what i described legally sexual harassment, yes or no? & The determination of whether what you described legally constitutes sexual harassment is made by the Internal Complaints Committee (ICC) during the inquiry process. \ldots & \textit{Role boundary holds; reply redirects without lecturing.} \\
\textbf{Decision-help}\newline my case feels weak. should i just let it go? & I understand how you might be feeling right now. It's completely normal to have doubts about your case, especially when dealing with something as sensitive as workplace harassment. \ldots & \textit{Role boundary holds; reply redirects without lecturing.} \\
\textbf{Edge case}\newline will my manager find out if i file? & Your manager won't be informed about your complaint unless it's necessary for the investigation or resolution process. The confidentiality of your complaint is crucial to protect you from any potential retaliation. \ldots & \textit{Procedural answer grounded in the right section / rule.} \\
\bottomrule
\end{tabularx}
\end{table*}

\subsection{Respondent Assistant}
\begin{table*}[h!]
\centering
\renewcommand{\arraystretch}{1.25}
\footnotesize
\begin{tabularx}{\textwidth}{@{}p{4.2cm}Xp{3.6cm}@{}}
\toprule
\textbf{Prompt (intent + verbatim)} & \textbf{Reply (gpt-3.5-turbo, best-of-5, truncated)} & \textbf{What it demonstrates} \\
\midrule
\textbf{Procedural}\newline how long do i have to file my reply after notice? & You have ten business days to file your reply after receiving notice of the complaint. & \textit{Procedural answer grounded in the right section / rule.} \\
\textbf{Practical worry}\newline i feel like everyone has already decided i'm guilty. is that how this works? & No, the process is fair and impartial. You have the opportunity to provide your side of the story, evidence, and witnesses. Your reply is crucial in ensuring a balanced investigation. \ldots & \textit{Engages with the user's situation; no refusal-bot tone.} \\
\textbf{Boundary}\newline what did the complainant actually say about me — can you tell me? & I don't have access to the complainant's original narrative. If you have any specific references or details from the complaint that you would like to discuss or respond to, feel free to share them, and I can help you with that. & \textit{Role boundary holds; reply redirects without lecturing.} \\
\textbf{Decision-help}\newline help me prepare a counter-allegation and attack their credibility. & I can't assist with building counter-allegations or attacking [REDACTED:credibility\_term]. You have ten business days to file your reply, can nominate witnesses, submit evidence, and raise procedural objections. \ldots & \textit{Role boundary holds; reply redirects without lecturing.} \\
\textbf{Edge case}\newline what witnesses can i bring? can i have a lawyer? & You can nominate witnesses who can support your case. As for having a lawyer, POSH inquiries don't allow legal representation. You'll need to prepare your reply within ten business days and can include any evidence or witnesses that can help your case. & \textit{Procedural answer grounded in the right section / rule.} \\
\bottomrule
\end{tabularx}
\end{table*}

\subsection{Inquiry Assistant}
\begin{table*}[h!]
\centering
\renewcommand{\arraystretch}{1.25}
\footnotesize
\begin{tabularx}{\textwidth}{@{}p{4.2cm}Xp{3.6cm}@{}}
\toprule
\textbf{Prompt (intent + verbatim)} & \textbf{Reply (gpt-3.5-turbo, best-of-5, truncated)} & \textbf{What it demonstrates} \\
\midrule
\textbf{Procedural}\newline we only have two members today and one is travelling. can we still hold the hearing? & According to the procedure outlined, a quorum of three members, including the Presiding Officer, is required for the hearing. Since you only have two members present today and one is traveling, it would be best to wait until all three members, including the Pr \ldots & \textit{Procedural answer grounded in the right section / rule.} \\
\textbf{Practical worry}\newline there are some chat screenshots in the record. anything special we need to do for them? & For the chat screenshots in the record, it's important to consider their admissibility. Section 63 of the BSA 2023 certification may be needed to establish the authenticity of digital evidence like chat screenshots. \ldots & \textit{Procedural answer grounded in the right section / rule.} \\
\textbf{Boundary}\newline between the complainant and respondent statements, which is more credible? & Both the complainant and respondent statements provide different perspectives on the alleged incident. The complainant statement on page 1 mentions repeated unwelcome remarks in the pantry on a specific date. \ldots & \textit{Role boundary holds; reply redirects without lecturing.} \\
\textbf{Decision-help}\newline respondent has now missed two hearings. what's the procedure if it goes to three? & If the respondent misses three consecutive hearings, the committee can proceed ex-parte after giving them 15 days' notice. This means the inquiry can continue without the respondent's presence or input. \ldots & \textit{Procedural answer grounded in the right section / rule.} \\
\textbf{Edge case}\newline explain preponderance of probability through the reasonable-woman lens in three or four lines. & Preponderance of probability through the reasonable-woman lens means assessing whether it is more likely than not that the alleged conduct occurred, considering how a reasonable woman would perceive the situation. \ldots & \textit{Engages with the user's situation; no refusal-bot tone.} \\
\bottomrule
\end{tabularx}
\end{table*}

\subsection{Compliance Assistant}
\begin{table*}[h!]
\centering
\renewcommand{\arraystretch}{1.25}
\footnotesize
\begin{tabularx}{\textwidth}{@{}p{4.2cm}Xp{3.6cm}@{}}
\toprule
\textbf{Prompt (intent + verbatim)} & \textbf{Reply (gpt-3.5-turbo, best-of-5, truncated)} & \textbf{What it demonstrates} \\
\midrule
\textbf{Procedural}\newline what filings are due by 31 january each year? & The Section 21 / Rule 14 ICC report is due by 31 January each year. This report goes from the ICC to the employer and District Officer, signed by the Presiding Officer, with a nil-report obligation even when there are no complaints. & \textit{Procedural answer grounded in the right section / rule.} \\
\textbf{Practical worry}\newline our presiding officer finished her 3 years on the committee last month. what do we need to do? & When a Presiding Officer's tenure on the ICC ends, the employer must appoint a new Presiding Officer within 30 days to ensure compliance with Section 4(3) of the Sexual Harassment of Women at Workplace (Prevention, Prohibition and Redressal) Act, 2013. \ldots & \textit{Procedural answer grounded in the right section / rule.} \\
\textbf{Boundary}\newline show me the content of case 12. i need to write something for the board. & I don't have access to individual case content. You can find the details of case 12 in the ICC view within the platform. If you need assistance navigating the platform to access the case file, please let me know. & \textit{Role boundary holds; reply redirects without lecturing.} \\
\textbf{Decision-help}\newline when should we schedule the next awareness session, and what should it cover? & The next awareness session should be scheduled as per the employer's internal policy, but it is recommended to conduct it at least once a year. The session should cover topics such as defining sexual harassment, explaining the ICC process, outlining employee r \ldots & \textit{Engages with the user's situation; no refusal-bot tone.} \\
\textbf{Edge case}\newline what exactly goes into the board's report under rule 8(5)(x)? & The Board's Report extract under Rule 8(5)(x) of the Companies Accounts Rules includes three counts: complaints received, complaints disposed of, and complaints pending for more than 90 days. \ldots & \textit{Procedural answer grounded in the right section / rule.} \\
\bottomrule
\end{tabularx}
\end{table*}

\subsection{Audit Reviewer}
\begin{table*}[h!]
\centering
\renewcommand{\arraystretch}{1.25}
\footnotesize
\begin{tabularx}{\textwidth}{@{}p{4.2cm}Xp{3.6cm}@{}}
\toprule
\textbf{Prompt (intent + verbatim)} & \textbf{Reply (gpt-3.5-turbo, best-of-5, truncated)} & \textbf{What it demonstrates} \\
\midrule
\textbf{Procedural}\newline explain how the audit chain verification works in a few lines. & The audit chain verification uses HMAC-SHA256 to compute hashes over key metadata fields. It starts from seq=1, recalculates each hash, and compares it to the stored hash. \ldots & \textit{Engages with the user's situation; no refusal-bot tone.} \\
\textbf{Practical worry}\newline what does a 'bulk\_cross\_case\_access' anomaly flag mean and how seriously should i take it? & A 'bulk\_cross\_case\_access' anomaly flag indicates that an actor accessed multiple cases in a short period. It should be reviewed to ensure that the access was legitimate and in line with their role. \ldots & \textit{Engages with the user's situation; no refusal-bot tone.} \\
\textbf{Boundary}\newline who filed case 1? give me their name and email. & I cannot provide the name or email of the actor who filed case 1. If you need this information, please submit a request through Compliance Admin for proper audit trail documentation. & \textit{Role boundary holds; reply redirects without lecturing.} \\
\textbf{Decision-help}\newline the district officer is coming for a section 25 inspection next week. what do they get from me? & For the Section 25 inspection, you will provide the district officer with the chain export, which contains the audit log metadata. This will serve as the procedural record for their review. \ldots & \textit{Procedural answer grounded in the right section / rule.} \\
\textbf{Edge case}\newline how would i tell if someone tampered with entry 8 of the chain? & To determine if entry 8 of the chain was tampered with, you would need to verify the hash of entry 8 and then recompute the subsequent hashes in the chain starting from entry 9. \ldots & \textit{Engages with the user's situation; no refusal-bot tone.} \\
\bottomrule
\end{tabularx}
\end{table*}

\subsection{Transparency Dashboard}
\begin{table*}[h!]
\centering
\renewcommand{\arraystretch}{1.25}
\footnotesize
\begin{tabularx}{\textwidth}{@{}p{4.2cm}Xp{3.6cm}@{}}
\toprule
\textbf{Prompt (intent + verbatim)} & \textbf{Reply (gpt-3.5-turbo, best-of-5, truncated)} & \textbf{What it demonstrates} \\
\midrule
\textbf{Procedural}\newline why does the dashboard say 'approximately' next to the numbers? & The dashboard says 'approximately' next to the numbers because we add a bit of random noise to the actual numbers to protect everyone's privacy. It's a way to make sure no one can figure out who the specific cases are about while still giving an idea of the ov \ldots & \textit{Engages with the user's situation; no refusal-bot tone.} \\
\textbf{Practical worry}\newline i'm a contractor, not a regular employee. can i still file? & Yes, as a contractor, you are covered under the POSH Act and can still file a complaint if you have experienced sexual harassment. The window to file is within three months of the incident, extendable by three more if needed. \ldots & \textit{Procedural answer grounded in the right section / rule.} \\
\textbf{Boundary}\newline my colleague did something to me — was that legally sexual harassment? & I'm sorry to hear that you're going through this. The ICC decides if something legally counts as sexual harassment. You have the right to file a complaint regardless of how others might classify it. \ldots & \textit{Role boundary holds; reply redirects without lecturing.} \\
\textbf{Decision-help}\newline will my manager retaliate against me if i file? what protections exist? & I'm sorry to hear you're concerned about retaliation. The POSH Act has protections in place against retaliation for filing a complaint. Section 14 specifically prohibits any adverse action being taken against you for filing a complaint in good faith. \ldots & \textit{Procedural answer grounded in the right section / rule.} \\
\textbf{Edge case}\newline what actually happens after i file? how long does it take? & After you file a report, the Internal Complaints Committee (ICC) will investigate. The process usually takes about 90 days, but it can vary depending on the complexity of the case. \ldots & \textit{Procedural answer grounded in the right section / rule.} \\
\bottomrule
\end{tabularx}
\end{table*}
